% texte.tex — f165r, vv. 208–251 (conversion initiale des travaux antérieurs)
\vers[n=209]{ΕΥ.}{Euripide}{ὦ τρισκακοδαίμον, ὡς ἀπόλωλ᾽ Εὐριπίδης.}{Ô triple malheureux, comment périt Euripide.}
\vers[n=210]{ΜΝ.}{Mnésiloque}{ὦ φίλτατ᾽ ὦ κηδεστὰ μὴ σαυτὸν προδῷς.}{Mon ami, mon gendre, n'abandonne point.}
\vers[n=211]{ΕΥ.}{Euripide}{πῶς οὖν ποιήσω δῆτα;}{Que puis-je faire ?}
\vers[suite]{ΜΝ.}{Mnésiloque}{τοῦτον μὲν μακρὰ}{Envoie-le se faire}
\vers[n=212]{}{}{κλάειν κέλευ᾽, ἐμοὶ δ᾽ ὅ τι βούλει χρῶ λαβών.}{pendre, lui, et dispose de moi comme tu voudras.}
\vers[n=213]{ΕΥ.}{Euripide}{ἄγε νυν ἐπειδὴ σαυτὸν ἐπιδίδως ἐμοί,}{Voyons, puisque tu te remets entre mes mains,}
\vers[n=214]{}{}{ἀπόδυθι τουτὶ θοἰμάτιον.}{ôte ce manteau.}
\vers[suite]{ΜΝ.}{Mnésiloque}{καὶ δὴ χαμαί.}{Le voilà par terre.}
\vers[n=215]{}{}{ἀτὰρ τί μέλλεις δρᾶν μ᾽;}{Mais que penses-tu faire de moi ?}
\vers[suite]{ΕΥ.}{Euripide}{ἀποξυρεῖν ταδί,}{Te raser cela}
\vers[n=216]{}{}{τὰ κάτω δ᾽ ἀφεύειν.}{et te flamber par dessous.}
\vers[suite]{ΜΝ.}{Mnésiloque}{ἀλλὰ πρᾶττ᾽, εἴ σοι δοκεῖ.}{Eh bien, fais à ta guise ;}
\vers[n=217]{}{}{ἢ μὴ διδόναι γ' ἑαυτὸν ὤφελόν ποτε.}{autrement je n'aurais jamais dû me mettre à ta merci.}
\vers[n=218]{ΕΥ.}{Euripide}{Ἀγάθων σὺ μέντοι ξυροφορεῖς ἑκάστοτε,}{Agathon, pour sûr, tu portes toujours sur toi des rasoirs.}
\vers[n=219]{}{}{χρῆσόν τι νῦν ἡμῖν ξυρόν.}{Prête-moi donc un peu un rasoir.}
\vers[suite]{ΑΓ.}{Agathon}{αὐτὸς λάμβανε}{Prends-le toi-même,}
\vers[n=220]{}{}{ἐντεῦθεν ἐκ τῆς ξυροδόκης.}{là, dans l'étui.}
\vers[suite]{ΕΥ.}{Euripide}{γενναῖος εἶ.}{Tu es galant homme.}
\vers[n=221]{}{}{κάθιξε· φύσα τὴν γνάθον τὴν δεξιάν.}{\didasc{À Mnésiloque.} Assieds-toi. Enfle la joue droite. \didasc{Il le rase.}}
\vers[n=222]{ΜΝ.}{Mnésiloque}{οἴμοι.}{Aïe !}
\vers[suite]{ΕΥ.}{Euripide}{τί κέκραγας; ἐμβαλῶ σοι πάτταλον,}{Qu'as-tu à crier ? Je t'enfonce un bâillon}
\vers[n=223]{}{}{ἢν μὴ σιωπᾷς.}{si tu ne te tais.}
\vers[suite]{ΜΝ.}{Mnésiloque}{ἀτταταῖ ἰατταταῖ.}{Aïe, aïe, aïe ! aie, aïe, aïe !}
\vers[n=224]{ΕΥ.}{Euripide}{οὗτος σὺ ποῖ θεῖς;}{Holà ! où cours-tu ?}
\vers[suite]{ΜΝ.}{Mnésiloque}{εἰς τὸ τῶν σεμνῶν θεῶν·}{Au temple des augustes}
\vers[n=225]{}{}{οὐ γὰρ μὰ τὴν Δήμητρά γ᾽ ἐνταυθοῖ μενῶ}{déesses. Car, par Déméter, je ne resterai pas là à me faire}
\vers[n=226]{}{}{τεμνόμενος.}{taillader.}
\vers[suite]{ΕΥ.}{Euripide}{οὔκουν καταγέλαστος δῆτ᾽ ἔσει}{Mais tu seras grotesque}
\vers[n=227]{}{}{τὴν ἡμίκραιραν τὴν ἑτέραν ψιλὴν ἔχων;}{avec ta moitié de figure dégarnie de poil.}
\vers[n=228]{ΜΝ.}{Mnésiloque}{ὀλίγον μέλει μοι.}{Peu m'en chaut.}
\vers[suite]{ΕΥ.}{Euripide}{μηδαμῶς πρὸς τῶν θεῶν}{Au nom des dieux,}
\vers[n=229]{}{}{προδῷς με· χώρει δεῦρο.}{ne m'abandonne pas. Arrive ici.}
\vers[suite]{ΜΝ.}{Mnésiloque}{κακοδαίμων ἐγώ.}{Misère de moi !}
\vers[n=230]{ΕΥ.}{Euripide}{ἔχ᾽ ἀτρέμα σαυτὸν κἀνάκυπτε· ποῖ στρέφει;}{Tiens-toi coi et lève la tête. Où te tournes-tu?}
\vers[n=231]{ΜΝ.}{Mnésiloque}{μυμῦ.}{Mumû.}
\vers[suite]{ΕΥ.}{Euripide}{τί μύζεις; πάντα πεποίηται καλῶς.}{Pourquoi mumû ? Tout est fini, et selon les règles.}
\vers[n=232]{ΜΝ.}{Mnésiloque}{οἴμοι κακοδαίμων, ψιλὸς αὖ στρατεύσομαι.}{Ah malheur ! Je vais resservir dans l'infanterie légère.}
\vers[n=233]{ΕΥ.}{Euripide}{μὴ φροντίσῃ· ὡς εὐπρεπὴς φανεῖ πάνυ.}{Ne t'inquiète pas. Tu auras un air tout avenant.}
\vers[n=234]{}{}{βούλει θεάσασθ᾽ αὑτόν;}{Veux-tu te mirer ? \didasc{Il lui tend un miroir.}}
\vers[suite]{ΜΝ.}{Mnésiloque}{εἰ δοκεῖ, φέρε.}{Si c'est ton plaisir. Passe.}
\vers[n=235]{ΕΥ.}{Euripide}{ὁρᾷς σεαυτόν;}{Te vois-tu ?}
\vers[suite]{ΜΝ.}{Mnésiloque}{οὐ μὰ Δί᾽ ἀλλὰ Κλεισθένη.}{Non, pas Zeus. Je vois Clisthène.}
\vers[n=236]{ΕΥ.}{Euripide}{ἀνίστασ᾽, ἵν᾽ ἀφεύσω σε, κἀγκύψας ἔχε.}{Lève-toi, que je te flambe. Tiens-toi penché.}
\vers[n=237]{ΜΝ.}{Mnésiloque}{οἴμοι κακοδαίμων δελφάκιον γενήσομαι.}{Ah malheur ! Je vais être cochon de lait.}
\vers[n=238]{ΕΥ.}{Euripide}{ἐνεγκάτω τις ἔνδοθεν δᾷδ᾽ ἢ λύχνον.}{Qu'on m'apporte de dedans une torche ou une lampe.}
\vers[n=239]{}{}{ἐπίκυπτε· τὴν κέρκον φυλάττου νυν ἄκραν.}{Courbe-toi, et gare au bout de la queue.}
\vers[n=240]{ΜΝ.}{Mnésiloque}{ἐμοὶ μελήσει νὴ Δία, πλήν γ᾽ ὅτι κάομαι.}{J'y veillerai, par Zeus... N'empêche que je brûle.}
\vers[n=241]{}{}{οἴμοι τάλας. ὕδωρ ὕδωρ ὦ γείτονες}{Ah misère ! de l'eau, de l'eau, voisins,}
\vers[n=242]{}{}{πρὶν ἀντιλαβέσθαι πρωκτὸν ἓτερον τῆς φλογός.}{avant que la fesse d'à côté ne s'enflamme.}
\vers[n=243]{ΕΥ.}{Euripide}{θάρρει.}{Du sang froid.}
\vers[suite]{ΜΝ.}{Mnésiloque}{τί θαρρῶ καταπεπυρπολημένος;}{Du sang froid ? quand je suis dévoré par les flammes !..}
\vers[n=244]{ΕΥ.}{Euripide}{ἀλλ᾽ οὐκ ἔτ᾽ οἰδὲν πρᾶγμά σοι· τὰ πλεῖστα γὰρ}{Mais tu n'as plus à t'inquiéter. Car te voilà dépêtré}
\vers[n=245]{}{}{ἀποπεπόνηκας.}{du principal.}
\vers[suite]{ΜΝ.}{Mnésiloque}{φεῦ ἰοὺ τῆς ἀσβόλου.}{Phu ! Ah! cette suie !}
\vers[n=246]{}{}{αἰθὸς γεγένημαι πάντα τὰ περὶ τὴν τράμιν.}{Je suis tout feu dans l'entrefesson.}
\vers[n=247]{ΕΥ.}{Euripide}{μὴ φροντίσῃς· ἕτερος γὰρ αὐτὰ σπογγιεῖ.}{Ne t'inquiète pas. Le voisin d'à côté t'épongera cela.}
\vers[n=248]{ΜΝ.}{Mnésiloque}{οἰμώξετ᾽ ἆρ᾽ εἴ τις τὸν ἐμὸν πρωκτὸν πλυνεῖ.}{Il en coûtera à celui qui s'attaquera à mon derrière.}
\vers[n=249]{ΕΥ.}{Euripide}{Ἀγάθων, ἐπειδὴ σαυτὸν ἐπιδοῦναι φθονεῖς,}{Agathon, puisque tu refuses de te dévouer à moi,}
\vers[n=250]{}{}{ἀλλ᾽ ἱμάτιον γοῦν χρῆσον ἡμῖν τουτῳὶ}{prête-nous du moins un vêtement}
\vers[n=251]{}{}{καὶ στρόφιον· οὐ γὰρ ταῦτά γ᾽ ὡς οὐκ ἔστ᾽ ἐρεῖς.}{et un strophion. Tu ne nieras pas que tu en possèdes..}
